\documentclass [a4paper,11pt,titlepage] {article}
\usepackage{graphicx}
\usepackage{pdfpages}
\usepackage{fancybox}
\usepackage[francais]{babel}
\usepackage[utf8]{inputenc}
% \usepackage[T1]{fontenc}
\usepackage{amsmath,amsfonts,amssymb}
\usepackage{fancyhdr}
\usepackage{stackrel}
\usepackage{xspace}
\usepackage{url}
\usepackage{titling}
\usepackage{listings}
\usepackage{color}
\usepackage{array}
\usepackage{hyperref}
\usepackage{makecell}
\newcolumntype{x}[1]{>{\centering\arraybackslash}p{#1}}
\usepackage{tikz}
\newcommand\diag[4]{%
  \multicolumn{1}{p{#2}|}{\hskip-\tabcolsep
  $\vcenter{\begin{tikzpicture}[baseline=0,anchor=south west,inner sep=#1]
  \path[use as bounding box] (0,0) rectangle (#2+2\tabcolsep,\baselineskip);
  \node[minimum width={#2+2\tabcolsep},minimum height=\baselineskip+\extrarowheight] (box) {};
  \draw (box.north west) -- (box.south east);
  \node[anchor=south west] at (box.south west) {#3};
  \node[anchor=north east] at (box.north east) {#4};
 \end{tikzpicture}}$\hskip-\tabcolsep}}

\definecolor{dkgreen}{rgb}{0,0.6,0}
\definecolor{gray}{rgb}{0.5,0.5,0.5}
\definecolor{mauve}{rgb}{0.58,0,0.82}

\lstset{frame=tb,
  language=Java,
  aboveskip=3mm,
  belowskip=3mm,
  showstringspaces=false,
  columns=flexible,
  basicstyle={\small\ttfamily},
  numbers=none,
  numberstyle=\tiny\color{gray},
  keywordstyle=\color{blue},
  commentstyle=\color{dkgreen},
  stringstyle=\color{mauve},
  breaklines=true,
  breakatwhitespace=true,
  tabsize=3
}

\setlength{\parindent}{0pt}
\setlength{\parskip}{1ex}
\setlength{\textwidth}{17cm}
\setlength{\textheight}{24cm}
\setlength{\oddsidemargin}{-.7cm}
\setlength{\evensidemargin}{-.7cm}
\setlength{\topmargin}{-.5in}


\predate{
\begin{center}
}
\postdate{
\\
\vspace{1.5cm}
\includegraphics[scale=0.7]{imag.png}
\end{center}}


\title {{ {\huge Compte rendu du projet}} \\
``{\em Simulation Orientée-Objet de systèmes multiagents}'' }

\author {Equipe 32 \\
\\
DOUMA Nejmeddine\\
GOUTTEFARDE Léo\\
KACHER Ilyes}
\date{Lundi 16 Novembre 2015}
% \date{Lundi 16 Novembre 2015\endgraf\bigskip
% Equipe NN}

\lhead{Projet POO}
\rhead{Compte rendu}

\begin{document}
\pagestyle{fancy}
\maketitle

\begin{center}
\section* {Introduction }
\end{center}


L'objectif de ce TP était d'implémenter en \texttt{Java} les simulations
graphiques de quelques systèmes multiagents, l'un des points les plus importants étant d'aborder les aspects fondamentaux de la programmation orientée objet.

Nous avons implémenté trois simulateurs d'automates cellulaires : le jeu de la vie de Conway, le jeu de l'immigration, le modèle de ségrégation de Schelling, ainsi qu'un simulateur de Boids (système d'essaims auto-organisés).


$\newline$
\section {Choix conceptuels}
\subsection{Organisation en packages}
Afin de faciliter la navigation au sein des différentes classes du projet, nous les avons réparties en différents packages :
\begin{itemize}
\item \texttt{element} : classes des différents agents (\texttt{Boid}, \texttt{Lighter}, \texttt{Predator} et \texttt{Prey})
\item \texttt{event} : classes relatives aux événements
\item \texttt{group} : classes qui gèrent les groupes d'agents
\item \texttt{simulator} : classes simulator qui implémentent \texttt{Simulable}
\item \texttt{test} : classes des tests
\item \texttt{utility} : classes annexes (\texttt{PVector}, \texttt{Triangle}, \texttt{Type})
\end{itemize}

\subsection{Premier simulateur : Balls}
Pour réaliser ce simulateur nous avons suivi le sujet qui est assez guidé, et notamment géré les collisions. Cette étape a permis de nous familiariser avec la librairie \texttt{gui.jar} fournie et en particulier l'interface \texttt{Simulable}.

\subsection{Automates cellulaires}
Nous avons cherché à factoriser notre code tout en profitant des avantages de la programmation orientée objet en implémentant deux classes abstraites \texttt{Automaton} et \texttt{AutomatonSimulator}. Pour chaque système d'automate cellulaire simulé il y a une classe concrète qui hérite de la classe \texttt{Automaton} et implémente les règles du système, ainsi qu'une autre pour la simuler en héritant de \texttt{AutomatonSimulator}.
Nous avons remarqué des similarités entre le jeu de l'immigration et le modèle du Schelling, ainsi, afin de factoriser davantage, nous avons défini une autre classe abstraite \texttt{ExtendedAutomaton} qui étend \texttt{Automaton} et est héritée par l'implémentation de ces deux systèmes.

\begin{center}
\begin{tabular}{|l|c|c|c|}
  \hline
    Système & Implémentation & Simulateur \\
  \hline
   Jeux de la vie de Conway & \texttt{Life} & \texttt{LifeSimulator}\\
  \hline
  Jeux de l'immigration & \texttt{Immigration} & \texttt{ImmigrationSimulator}\\
  \hline
  Modèle de Schelling & \texttt{Schelling} & \texttt{SchellingSimulator}\\
  \hline
 
\end{tabular}
\end{center}

Nous avons également généralisé l'utilisation de \texttt{EventManager} au calcul des générations suivantes pour tous les systèmes d'automate cellulaire simulés en définissant \texttt{AutomatonEvent} qui hérite de la classe abstraite \texttt{Event}.


\subsection{Mouvements en essaims auto-organisés : les Boids}
\subsubsection {Le système de {\em boids}}
Nous avons presque utilisé la même architecture que pour le simulateur des balles. La classe abstraite \texttt{Boid} implémente les règles communes d'un agent boid, la classe \texttt{Boids} implémente un environnement de plusieurs {\em boids} tandis que le simulateur du système se trouve dans la classe \texttt{BoidsSimulator}. 


\subsubsection {Les types de {\em boids}}

Afin de compléter le comportement d'un boid nous avons ajouté les classes \texttt{Prey}, \texttt{Predator} et \texttt{Lighter}. Ces classes héritent de \texttt{Boid} et implémentent les règles nécessaires afin d'obtenir un boid proie ou prédateur. Le boid \texttt{Lighter} est un boid rond et brillant dont la couleur change à intervalles de temps aléatoires (événements différés aléatoires).


\subsubsection {Le système de Proies / Prédateurs}

Pour répondre à la simulation du système de proies et de prédateurs nous avons ajouté les règles suivantes :
\begin{itemize}
\item Un prédateur chasse la première proie de son champ de vision à sa vitesse maximale
\item Un prédateur dévore une proie suffisamment proche
\item Une proie s'enfuit à sa vitesse maximale dès qu'un prédateur entre dans son champ de vision
\end{itemize}

\section {Tests}

\subsection {Exécution}
Toutes les classes de test sont dans le package \texttt{test}, elles sont réparties selont le tableau suivant :

\begin{center}
\begin{tabular}{|l|c|c|}
  \hline
    Classe & Système simulé & Commande pour l'exécution \\
  \hline
  \texttt{TestGUI} & Test de l'interface graphique & \texttt{make exeGUI}\\
  \hline
   \texttt{TestBalls} & Test demandé à la première question & \texttt{make exeBalls}\\
  \hline
  \texttt{TestBallsSimulator} & Balls & \texttt{make exeBSim}\\
  \hline
  \texttt{TestLifeSimulator} & Jeux de la vie de Conway & \texttt{make exeLifeSim}\\
  \hline
  \texttt{TestImmigrationSimulator} & Jeu de l'immigration & \texttt{make exeImmSim}\\
  \hline
  \texttt{TestSchellingSimulator} & Modèle de Shelling & \texttt{make exeSSim}\\
  \hline
  \texttt{TestEventManager} & Test de l'event manager & \texttt{make exeEvents}\\
  \hline
  \texttt{TestBoidsSimulator} & Boids simples & \texttt{make exeBoids}\\
  \hline
  \texttt{TestPreyPredatorSimulator} & Boids proies/prédateurs  & \texttt{make exePred}\\
  \hline

\end{tabular}
\end{center}

\subsection {Résultats}
\subsubsection {Jeu de la vie de Conway}
Pour tester ce système, nous avons implémenté certains \textit{pattern} classiques comme \textit{toad}, \textit{pulsar} ou \textit{blinker} obtenus à l'adresse suivante : \url{https://en.wikipedia.org/wiki/Conway's_Game_of_Life#Examples_of_patterns}.
\newline
Nous avons bien obtenu les résultats attendus avec le bon nombre d'itérations.

\subsubsection {Jeux de l'immigration}
Nous avons pu constater l'effet des règles du jeux sur les cellules, en particulier la manière dont un état se propage sur la grille.

\subsubsection {Modèle de Schelling}
Le tableau suivant contient le nombre d'itérations nécessaires avant d'obtenir une ségrégation en fonction du seuil \textit{K} et du nombre de couleurs:
\begin{center}
\begin{tabular}{|x{5cm}|x{2cm}|x{2cm}|x{2cm}|x{2cm}|x{2cm}|}\hline
\diag{.1em}{5cm}{\texttt{NbColors}}{\textit{K}} & 2 & 4 & 5 & 7 & 8\\ \hline
2 & 55 & 16 & 8 & 0 & 0\\ \hline
7 & +10000 $\infty$ & 409 & 50 & 11 & 0\\ \hline
14 & +10000 $\infty$ & 2363 & 237 & 17 & 0\\ \hline
49 & +10000 $\infty$ & +10000 $\sim$ & 2270 & 40 & 0\\ \hline

\end{tabular}
\end{center}
+10000 $\sim$ : Nous avons dépassé 10000 itérations en n'ayant qu'une ségrégation partielle.
\newline
+10000 $\infty$ : Nous avons dépassé 10000 itérations sans aucune ségrégation en cours de formation.
\newline
\newline
Nous constatons que le nombre d'itérations nécessaires avant d'obtenir une ségrégation est inversement proportionnel à \textit{K} et proportionnel à \texttt{NbColors}.

\subsubsection {Boids}
Sur une même grille, nous avons lancé deux groupes de {\em boids} différents: des {\em boids} triangulaires (les \texttt{Prey}) et des {\em boids} circulaires (les \texttt{Lighter}). Nous avons réussi à visualiser la cohabitation entre ces deux groupes différents au sein d'une même simulation.

\subsubsection {Système Proie / Prédateur}
Les \texttt{Predator} cherchent à attraper les \texttt{Prey} sans se laisser gêner par les \texttt{Lighter} qui partagent la même grille. On voit bien le rôle joué par l'event manager consistant à gérer la mise à jour des agents selon différentes fréquences.
\\

\begin{center}
\section* {Conclusion}
\end{center}

Ce projet nous a introduit aux problématiques de la programmation orientée objet ainsi qu'à la manière d'en tirer profit pour produire un code à la fois factorisé et réutilisable.
Il nous a aidé non seulement à mieux comprendre les notions de cours mais aussi à découvrir le monde des simulateurs multi-agents.

C'est un projet intéressant qui nous a permis de profiter d'une interface visuelle et de créer des agents amusants comme le \texttt{Lighter}.

\end{document}




% 4 pages max
% Explique / justifie choix conception / bonne utilisation des classes et méthodes
% Peut également décrire principaux tests effectués et résultats obtenus

